\expandafter\def\csname HistoricalNoteText8\endcsname{Als ich vor einiger Zeit das Glück hatte, Professor \textsc{Gauss} persönlich kennenzulernen, sah ich unter seinen Papieren den folgenden Aufsatz. Er war bereits vor mehreren Jahren entworfen worden und noch nirgends veröffentlicht. Er enthält die frühere Methode des Verfassers zur Bestimmung von Planetenbahnen. Bei der flüchtigen Durchsicht dieser Übersicht war ich bald überzeugt, dass sich die hier entwickelte Methode zur ersten näherungsweisen Bestimmung zweier Abstände des Planeten von der Erde wesentlich von der Methode unterscheidet, die der Verfasser jetzt in seinem größeren Werk öffentlich dargelegt hat. Deshalb bat ich ihn um Erlaubnis, den Aufsatz zu veröffentlichen. Ich ging davon aus, dass alle Kenner daran interessiert sein müssen, die verschiedenen Wege zu erfahren, auf denen der Verfasser zu der vollständigen Lösung gelangt ist, von der wir unseren Lesern im vorigen Heft einen Überblick gegeben haben. Zunächst wollte ich den Aufsatz mit einigen Bemerkungen zum Vergleich der früheren und der späteren Methode des Verfassers begleiten. Um die Sache wirklich zu erläutern, hätten sie aber etwas ausführlich werden müssen und wären ohne Verweis auf das Werk selbst doch unklar geblieben. Deshalb schien es mir zweckmäßiger, den astronomischen Lesern dieser Zeitschrift den ganzen Aufsatz, der vor allem für Kenner bestimmt ist, die das Werk selbst zur Hand haben, ohne weiteren Zusatz so mitzuteilen, wie ihn der Verfasser vor sechs Jahren niedergeschrieben hat.\hfill\textsc{von Lindenau.}}
\clearpage\SourceLeaf{153}%
\pdfbookmark[1]{Kurzer Überblick über die Methoden zur Bestimmung der Bahnen der beiden neuen Hauptplaneten}{summarische1809}
\fancyhead[L]{\small METHODEN ZUR BAHNBESTIMMUNG}
\begin{center}
{\small KURZER ÜBERBLICK ÜBER DIE METHODEN ZUR BESTIMMUNG DER BAHNEN DER BEIDEN NEUEN HAUPTPLANETEN\HistoricalNote{8}.\par}
\medskip\hrule\medskip
{\small Monatliche Correspondenz zur Beförderung der Erd- und Himmels-Kunde,\par
herausgegeben vom Freiherrn \textsc{von Zach}. September 1809.\par}
\medskip\hrule
\end{center}\medskip
\begin{center}1.\end{center}
Die Bestimmung der Bahn eines Himmelskörpers aus einer kurzen Reihe von Beobachtungen, ohne eine Kreis- oder Parabelbahn vorauszusetzen, beruht auf zwei Forderungen. I. Man muss die Bahn finden können, die drei gegebenen vollständigen Beobachtungen genügt. II. Man muss die so gefundene Bahn so verbessern können, dass die Abweichungen der berechneten Positionen von sämtlichen Beobachtungen möglichst gering werden.

\SourceLeaf{154}%
Die einfachste Art, der II\textsuperscript{ten} Forderung zu genügen, scheint darin zu bestehen, sie auf die I\textsuperscript{ste} zurückzuführen. Zu den Zeiten $t$, $t'$, $t''$ usw. seien $m$, $m'$, $m''$ die beobachteten Positionen (mit je zwei Komponenten). Die nach den bekannten Elementen $e$ (die man als aus sechs Größen bestehend auffasst) berechneten Positionen seien $p$, $p'$, $p''$ usw. Schließlich seien die nach den (noch als unbestimmt anzusehenden) Elementen $f$ berechneten Positionen $q$, $q'$, $q''$ usw. Die Differenzen für die Elemente $e$ sind also
\[
p-m,\quad p'-m',\quad p''-m''\ \text{usw.}
\]
Die Differenzen für die Elemente $f$ lauten dagegen
\[
(p-m)+(q-p),\quad(p'-m')+(q'-p'),\quad(p''-m'')+(q''-p'')\ \text{usw.}
\]
Diese sollen nun möglichst klein werden und keine Regelmäßigkeit behalten.

Die Differenzen $q-p$, $q'-p'$ usw. sind Funktionen der Zeit, sofern man die Elemente $f$ als konstant ansieht. Da sie von der Natur der Sache her klein sein werden, darf man bei der kurzen Dauer der Beobachtungen annehmen, dass sie sich für die dazwischenliegenden Zeiten durch Interpolation hinreichend genau bestimmen lassen, wenn man sie für zwei äußere und eine mittlere Zeit als gegeben annimmt. Man bezeichne sie für diese drei Zeiten mit $x$, $y$, $z$ (mit je zwei Komponenten). Nach bekannten Grundsätzen der Interpolationstheorie haben sie dann die lineare Form $\alpha x+\beta y+\gamma z$, wobei die Koeffizienten $\alpha$, $\beta$, $\gamma$ von der Zeit abhängen. Diese Differenzen für die Elemente $f$ haben für die Zeiten $t$, $t'$, $t''$ usw. also die folgende Gestalt.
\[
\begin{gathered}
p-m+\alpha x+\beta y+\gamma z\\
p'-m'+\alpha'x+\beta'y+\gamma'z\\
p''-m''+\alpha''x+\beta''y+\gamma''z\ \text{usw.}
\end{gathered}
\]
Hier ist alles außer $x$, $y$, $z$ bekannt. Man kann dann leicht beurteilen, welche Werte für $x$, $y$, $z$ am zweckmäßigsten sind. \emph{Es lässt sich zwar ein ganz systematisches Rechenverfahren zur Bestimmung dieser Werte angeben.} Ein gewisses Gespür führt aber immer ebenso sicher.

Da also offenbar die II\textsuperscript{te} Forderung auf die erste zurückgeführt ist, sobald $x$, $y$, $z$ bestimmt sind, können wir uns auf die erste beschränken.

\SourceLeaf{155}%
\begin{center}\emph{Bestimmung der Bahn aus drei vollständigen Beobachtungen.}\end{center}
\begin{center}2.\end{center}
Es wäre zwar leicht, den Zusammenhang zwischen den sechs unbekannten und den gegebenen Größen in sechs Gleichungen darzustellen. Diese wären jedoch viel zu unhandlich, um überhaupt brauchbar zu sein. Daher muss man sich damit begnügen, schrittweise zu der Bahn zu gelangen, die die drei Beobachtungen genau wiedergibt. Offensichtlich müssen alle Verfahren, die dafür überhaupt geeignet sind, am Ende dasselbe Ergebnis liefern. Die Güte des Endergebnisses misst daher die Genauigkeit der zugrunde gelegten Beobachtungen. Über den Wert des Verfahrens sagt sie nichts aus. Der Wert des Verfahrens lässt sich nur danach beurteilen, wie viele Schritte nötig sind und wie leicht sie sich ausführen lassen. Ließen sich die drei Beobachtungen mit einem Verfahren nicht jederzeit nach Belieben genau wiedergeben, so wäre es schon zu viel gesagt, es als schlechteres Verfahren zu bezeichnen. Es wäre überhaupt kein Verfahren.

Die Untersuchung zerfällt also in zwei Teile, eine \emph{erste Näherung} und die \emph{Korrekturverfahren}. Die erste Näherung stützt sich auf gewisse nahezu zutreffende Beziehungen, die sich aus der Natur der Aufgabe ergeben. Sie weichen umso weniger ab, je näher die Beobachtungen beieinanderliegen, und sind, mathematisch gesprochen, für unendlich nahe Beobachtungen streng exakt. Je näher also die zugrunde gelegten Beobachtungen beieinanderliegen, desto geringer ist der Einfluss der Abweichung dieser Beziehungen von der Wahrheit. Das erleichtert die Korrekturverfahren oder macht sie ganz entbehrlich. Man muss jedoch bedenken, dass bei nahe beieinanderliegenden Beobachtungen kleine Beobachtungsfehler sehr stark, bisweilen gewaltig auf die Elemente wirken. Dadurch wird die spätere Verbesserung anhand der ganzen Beobachtungsreihe, die oben als zweite Forderung bezeichnet wurde, umso schwieriger. Allgemeine Regeln für die zweckmäßigste Auswahl der Beobachtungen lassen sich kaum angeben. Ihre Güte, ihre mehr oder weniger günstige Lage und die Art der Bahn selbst sind zu berücksichtigen. Bei Ceres ließen sich die äußersten, 41 Tage auseinanderliegenden Beobachtungen mit gutem Erfolg sofort für die erste Näherung verwenden, und die Korrekturrechnungen waren gar nicht mühsam. Auch bei der Berechnung der II\textsuperscript{ten} Pallasbahn ließ man die vorherigen Näherungen beiseite und wandte das Verfahren der ersten
\SourceLeaf{156}%
Näherung sofort erneut auf die Beobachtungen über 27 Tage an. Bei Bahnen, die einer Parabel näher kämen und bei denen die geozentrische Bewegung sehr schnell ist, würde man wohl lieber mit etwas kürzeren Zeitabständen beginnen. Ein durch Erfahrung geschultes Urteil ist hier der beste Wegweiser.

\begin{center}
\emph{Hauptschritte der ersten Näherung.}\par\medskip
Erster Schritt.\par\smallskip
\emph{Näherungsweise Bestimmung der Abstände von der Erde zu den beiden äußeren Beobachtungszeiten.}
\end{center}
\begin{center}3.\end{center}
Damit man bei den vielen Zeichen, die in der folgenden Untersuchung nötig sind, leichter den Überblick behält, werden entsprechende Größen für die Erde $P$ und den beobachteten Planeten $p$ durchweg mit denselben Buchstaben bezeichnet, für die Erde jedoch mit großen und für den Planeten mit kleinen. Kommt derselbe Buchstabe ohne Strich sowie mit einem und mit zwei Strichen vor, so ist anzunehmen, dass sich die zweite und dritte Form zur zweiten und dritten Beobachtung zu den Zeiten $\tau'$, $\tau''$ ebenso verhalten wie die erste Form zur Beobachtung zur Zeit $\tau$. Die Zeit $\tau'$ muss an sich nicht zwischen $\tau$ und $\tau''$ liegen. Die folgenden Vorschriften lassen sich aber am günstigsten anwenden, wenn $\tau'$ etwa in der Mitte zwischen $\tau$ und $\tau''$ liegt.

$S$ bezeichnet die als unbeweglich betrachtete Position der Sonne (im Raum). $p$ ist die Position des Planeten $p$ zur Zeit $\tau$. Daraus erklären sich $p'$, $p''$, $P$, $P'$, $P''$.

$\mathfrak{X}$, $\mathfrak{Y}$, $\mathfrak{Z}$ sind drei beliebige feste Ebenen, die einander im Mittelpunkt der Sonne senkrecht schneiden.

$x$, $y$, $z$ sind die senkrechten Abstände des Planeten $p$ von diesen drei Ebenen zur Zeit $\tau$. Daraus ergeben sich $x'$, $y'$, $z'$; $x''$, $y''$, $z''$; $X$, $Y$, $Z$; $X'$, $Y'$, $Z'$; $X''$, $Y''$, $Z''$.
\[
\left.\begin{aligned}
\xi&=x-X\\
\eta&=y-Y\\
\zeta&=z-Z
\end{aligned}\right|\quad\text{Hieraus }\xi',\ \eta',\ \zeta';\ \xi'',\ \eta'',\ \zeta''
\]

$\xi$, $\eta$, $\zeta$ sind also die senkrechten Abstände des Planeten $p$ von drei beweglichen Ebenen, die parallel zu $\mathfrak{X}$, $\mathfrak{Y}$, $\mathfrak{Z}$ durch $P$ gelegt werden.

\SourceLeaf{157}%
\[
\begin{array}{c|c|c|c|c|l}
r&&\text{des }p&&S&\text{alle positiv.}\\
R&\text{Abstände}&\text{des }P&\text{von}&S&\text{Hieraus }r'\ \text{usw.}\\
\rho&&\text{des }p&&P&
\end{array}
\]
\[
\begin{array}{c|c|c|c|l}
b&&Sp&&\mathfrak{Z}\\
B&\text{Winkel der Linie}&SP&\text{mit der Ebene}&\mathfrak{Z}\\
\beta&&Pp&&\text{die }\mathfrak{Z}\text{ parallel ist}
\end{array}
\]
\[
\left.\begin{aligned}
d&=r\cos b\\D&=R\cos B\\\delta&=\rho\cos\beta
\end{aligned}\right|\quad
\begin{gathered}\text{d.\ h. die auf die Ebene }\mathfrak{Z}\\\text{und die zu ihr parallele Ebene projizierten Abstände}\end{gathered}
\]
\[
\begin{array}{c|c|l}
l&&\mathfrak{Y}\\
L&\text{Winkel dieser Projektion mit der Ebene}&\mathfrak{Y}\\
\lambda&&\text{die zu }\mathfrak{Y}\text{ parallel ist.}
\end{array}
\]
Die Winkel $b$ und $l$ sind auf denselben Seiten von $\mathfrak{Z}$ und $\mathfrak{Y}$ als positiv anzusetzen, auf denen man $z$ und $y$ als positiv ansieht. Die Winkel $b$ kann man stets zwischen den Grenzen $-90^0$, $+90^0$ lassen (damit die $d$ usw. stets positiv bleiben). Die Winkel $l$ dagegen kann man von $0$ bis $360^0$ anwachsen lassen, und zwar so, dass man sie dort $\left\{\begin{smallmatrix}0\\180^0\end{smallmatrix}\right\}$ setzt, wo die $x$ $\left\{\begin{smallmatrix}\text{positiv}\\\text{negativ}\end{smallmatrix}\right\}$ sind. So erhält man
\[
\left.\begin{aligned}
x&=r\cos b\cos l=d\cos l\\
y&=r\cos b\sin l=d\sin l\\
z&=r\sin b\phantom{\cos l}=d\operatorname{tang}b
\end{aligned}\right|\quad
\begin{gathered}\text{und ähnliche Gleichungen für}\\x'\ \text{usw.}\quad X\ \text{usw.}\quad\xi\ \text{usw.}\end{gathered}
\]

Wir nehmen an, dass die Bahnen von $p$ und $P$ in Ebenen liegen, und sehen dabei von den äußeren Einwirkungen ab, die sie beeinflussen. Wir setzen die Länge von $p$ in seiner Bahn zur Zeit $\tau$ $=v$ (also $v'$, $v''$; $V$, $V'$, $V''$) und setzen $\frac12r'r''\sin(v''-v')=f$, $\frac12r''r\sin(v-v'')=f'$, $\frac12rr'\sin(v'-v)=f''$. Damit sind $f$, $-f'$, $f''$ die positiven Flächeninhalte der Dreiecke $p'Sp''$, $pSp''$, $pSp'$ (vorausgesetzt, dass $p$ rechtläufig ist und $\tau'$ zwischen $\tau$ und $\tau''$ liegt. Die Vorzeichen lassen sich in anderen Fällen ohne Schwierigkeit zuordnen). Entsprechend $F$, $F'$, $F''$. Mit $g$, $-g'$, $g''$, $G$, $-G'$, $G''$ bezeichnen wir die Flächeninhalte der Ausschnitte aus der ganzen Bahn, die diesen Dreiecken entsprechen. Wir setzen voraus, dass ihre Vorzeichen mit denen von $f$, $-f'$, $f''$, $F$, $-F'$, $F''$ übereinstimmen. Daher sind $g$, $g'$, $g''$ und $G$, $G'$, $G''$ zu den Zeitabständen $\tau''-\tau'$, $\tau-\tau''$, $\tau'-\tau$ proportional.

\SourceLeaf{158}%
Die Bahnen von $p$ und $P$ sind Kegelschnitte, deren große Halbachsen wir mit $a$, $A$ bezeichnen. Die Exzentrizität der Bahn von $p$ setzen wir $=e=\sin\varphi$ (für eine Ellipse). Daher ist $a\cos\varphi^2=k$ der halbe Parameter. Länge des Aphels von $p$ in seiner Bahn $=\pi$. Mittlere Länge $=m$ (also $m'$, $m''$, $M$, $M'$, $M''$). Die übrigen Zeichen werden im Verlauf der Untersuchung eingeführt.

\begin{center}4.\end{center}
Da $p$, $p'$, $p''$ und $S$ in einer Ebene liegen, folgt nach einem bekannten Lehrsatz
\[
0=xy'z''+x'y''z+x''yz'-xy''z'-x'yz''-x''y'z
\]
und daraus folgt, dass von den folgenden neun Größen
\[
\begin{array}{ccc}
y'z''-y''z'&y''z-yz''&yz'-y'z\\
z'x''-z''x'&z''x-zx''&zx'-z'x\\
x'y''-x''y'&x''y-xy''&xy'-x'y
\end{array}
\]
die drei oberen jeweils zu den drei mittleren und zu den drei unteren proportional sind. Man kann sich leicht von Folgendem überzeugen.

I. Auch $f$, $f'$, $f''$ sind zu eben diesen Größen proportional. Die drei oberen, mittleren und unteren Größen stellen nämlich jeweils nur die doppelten Flächeninhalte der Projektionen der Dreiecke mit den Flächeninhalten $f$, $f'$, $f''$ auf die Fundamentalebenen $\mathfrak{X}$, $\mathfrak{Y}$, $\mathfrak{Z}$ dar. Das Verhältnis dieser Größen zu den Flächeninhalten der Dreiecke ist also jeweils das des doppelten Kosinus der Neigung der Bahn von $p$ zur jeweiligen Ebene zur Einheit. (In einer vollständigen Abhandlung wären hier noch Bemerkungen zu den Vorzeichen nötig, die sich aber durch bloßes Rechnen leicht umgehen lassen.)

II. Multipliziert man die drei oberen, die drei mittleren oder die drei unteren Größen jeweils mit $x$, $x'$, $x''$, mit $y$, $y'$, $y''$ beziehungsweise mit $z$, $z'$, $z''$, so wird die Summe der Produkte $=0$. Daraus lässt sich leicht schließen
\[
\begin{array}{r|c|l}
0=fx+f'x'+f''x''&\text{Durch genau}&0=FX+F'X'+F''X''\\
0=fy+f'y'+f''y''&\text{entsprechende Schlüsse}&0=FY+F'Y'+F''Y''\\
0=fz+f'z'+f''z''&\text{erhält man}&0=FZ+F'Z'+F''Z''
\end{array}
\]
\SourceLeaf{159}%
Daraus lassen sich leicht die folgenden drei Gleichungen ableiten
\[
\begin{aligned}
&(F+F'')(f\xi+f'\xi'+f''\xi'')\\
&\qquad=(Ff''-F''f)(X-X'')+[F'(f+f'')-(F+F'')f']X'\\
&(F+F'')(f\eta+f'\eta'+f''\eta'')\\
&\qquad=(Ff''-F''f)(Y-Y'')+[F'(f+f'')-(F+F'')f']Y'\\
&(F+F'')(f\zeta+f'\zeta'+f''\zeta'')\\
&\qquad=(Ff''-F''f)(Z-Z'')+[F'(f+f'')-(F+F'')f']Z'
\end{aligned}
\]

Aus diesen drei Gleichungen leiten wir vier weitere wie folgt ab
\[
\begin{array}{c|c|c|c|l}
\text{zuerst mit}&\text{dann mit}&\text{dann mit}&\text{endlich mit}&\text{multiplizieren}\\
\eta\zeta''-\eta''\zeta&\eta Z'-Y'\zeta&\eta'Z'-Y'\zeta'&\eta''Z'-Y'\zeta''&\text{und die}\\
\zeta\xi''-\zeta''\xi&\zeta X'-Z'\xi&\zeta'X'-Z'\xi'&\zeta''X'-Z'\xi''&\text{Produkte}\\
\xi\eta''-\xi''\eta&\xi Y'-X'\eta&\xi'Y'-X'\eta'&\xi''Y'-X'\eta''&\text{addieren.}
\end{array}
\]

Der Übersicht halber bezeichnen wir die Produktsummen, die entstehen
\begin{center}\begingroup\fontsize{9}{11}\selectfont\setlength{\tabcolsep}{.8pt}
\begin{tabular}{c|c|c|c|c|c}
&&&&&\shortstack{nach\\Multiplikation\\dieser\\Faktoren mit}\\\hline
&$*$&$*$&$\delta'\delta D'[\pi'\pi P']$&$\delta''\delta D'[\pi''\pi P']$&$\xi,\quad\eta,\quad\zeta$\\
$[\pi\pi'\pi'']$&$\delta\delta'\delta''\times[\pi\pi'\pi'']$&$\delta\delta'D'[\pi\pi'P']$&$*$&$\delta''\delta'D'[\pi''\pi'P']$&$\xi',\quad\eta',\quad\zeta'$\\
&$*$&$\delta\delta''D'[\pi\pi''P']$&$\delta'\delta''D'[\pi'\pi''P']$&$*$&$\xi'',\quad\eta'',\quad\zeta''$\\
$[\pi P\pi'']$&$\delta D\delta''[\pi P\pi'']$&$\delta DD'[\pi PP']$&$\delta'DD'[\pi'PP']$&$\delta''DD'[\pi''PP']$&$X,\quad Y,\quad Z$\\
$[\pi P'\pi'']$&$\delta D'\delta''[\pi P'\pi'']$&$*$&$*$&$*$&$X',\quad Y',\quad Z'$\\
&$\delta D''\delta''[\pi P''\pi'']$&$\delta D''D'[\pi P''P']$&$\delta'D''D'[\pi'P''P']$&$\delta''D''D'[\pi''P''P']$&$X'',\quad Y'',\quad Z''$
\end{tabular}\endgroup\end{center}
Offenbar muss an den hier mit $*$ bezeichneten Stellen $0$ stehen, und alle durch eingeklammerte Zeichen bezeichneten Größen sind \emph{gegeben}. Es gilt nämlich
\[
\begin{aligned}\relax
[\pi\pi'\pi'']&=\operatorname{tang}\beta\sin(\lambda'-\lambda'')+\operatorname{tang}\beta'\sin(\lambda''-\lambda)+\operatorname{tang}\beta''\sin(\lambda-\lambda')\\
[\pi P\pi'']&=\operatorname{tang}\beta\sin(L-\lambda'')+\operatorname{tang}B\sin(\lambda''-\lambda)+\operatorname{tang}\beta''\sin(\lambda-L)\\
&\text{usw.}
\end{aligned}
\]

Man braucht nicht alle 16 Gleichungen hinzuschreiben, da man sie auf entsprechende Weise aus der ersten ableiten kann, indem man nur $\beta$ durch
\expandafter\def\csname HistoricalNoteText9\endcsname{Der Kürze halber erlaube ich mir diesen leicht verständlichen Ausdruck. Der positive Pol von $\mathfrak{X}$ liegt auf der Seite dieser Ebene, auf der die $x$ positiv gezählt werden, und so fort.}%
\SourceLeaf{160}%
$\beta'$, $\beta''$, $B$, $B'$, $B''$ und $\lambda$ durch $\lambda'$, $\lambda''$, $L$, $L'$, $L''$ ersetzt, wenn an der Stelle von $\pi$ jeweils $\pi'$, $\pi''$, $P$, $P'$, $P''$ steht, und so fort. Zugleich sieht man, dass sich die 16 Größen auf 12 reduzieren, denn es gilt
\[
\begin{aligned}
+[\pi P'\pi'']&=-[\pi\pi''P']=+[\pi''\pi P']\\
[\pi\pi'P']&=-[\pi'\pi P']\\
[\pi''\pi'P']&=-[\pi'\pi''P']
\end{aligned}
\]

Ferner erkennt man leicht, dass der Ausdruck $[\pi\pi'\pi'']$, multipliziert mit dem Produkt der drei Kosinus der darin vorkommenden Breiten, das Sechsfache des Volumens einer Pyramide ist. Ihre Spitze liegt im Mittelpunkt einer Kugel mit dem Radius $1$, und die drei Eckpunkte ihrer Grundfläche liegen so auf der Kugeloberfläche, dass sie den drei geozentrischen Positionen von $p$ entsprechen. Jenes Zeichen gibt das Sechsfache des Volumens mit positivem oder negativem Vorzeichen an, je nachdem, ob die drei geozentrischen Positionen auf der Sphäre in entgegengesetzter oder in gleicher Reihenfolge (sens) liegen wie die positiven\HistoricalNote{9} Pole der Ebenen $\mathfrak{X}$, $\mathfrak{Y}$, $\mathfrak{Z}$ in dieser Reihenfolge. Die übrigen Zeichen drücken ganz entsprechende Größen aus.

So ergeben sich die folgenden vier Gleichungen:
\[
\begin{aligned}
1)\quad &(F+F'')f'\delta'[\pi\pi'\pi'']\\
&=(Ff''-F''f)(D[\pi P\pi'']-D''[\pi P''\pi''])\\
&\qquad+(F'(f+f'')-(F+F'')f')D'[\pi P'\pi'']\\
2)\quad &(F+F'')(f'\delta'[\pi\pi'P']+f''\delta''[\pi\pi''P'])\\
&=(Ff''-F''f)(D[\pi PP']-D''[\pi P''P'])\\
3)\quad &(F+F'')(f\delta[\pi'\pi P']+f''\delta''[\pi'\pi''P'])\\
&=(Ff''-F''f)(D[\pi'PP']-D''[\pi'P''P'])\\
4)\quad &(F+F'')(f\delta[\pi''\pi P']+f'\delta'[\pi''\pi'P'])\\
&=(Ff''-F''f)(D[\pi''PP']-D''[\pi''P''P'])
\end{aligned}
\]

\begin{center}5.\end{center}
Wir wollen nun diese vier Gleichungen, die streng richtig sind, näher betrachten, um darauf die erste Näherung zu gründen. Betrachten wir die Zeitabstände zwischen den Beobachtungen als unendlich kleine Größen erster Ordnung, so sind $f$, $f'$, $f''$, $G$, $G'$, $G''$ und alle eingeklammerten Größen ebenfalls von erster Ordnung, mit Ausnahme von $[\pi\pi'\pi'']$, das von dritter Ordnung ist. Ich übergehe die Beweise und die naheliegenden Bemerkungen zu besonderen Ausnah%
\SourceLeaf{161}%
men. Müsste man die gegenseitige Neigung der Bahnen von $p$ und $P$ als eine Größe erster Ordnung betrachten, so wären alle eingeklammerten Größen von einer Ordnung höher. Ferner ist
\[
Ff''-F''f=\frac FG.\frac{f''}{g''}Gg''-\frac{F''}{G''}.\frac fg G''g
\]
oder (weil $Gg''=G''g$)
\[
=\left(\frac FG.\frac{f''}{g''}-\frac{F''}{G''}.\frac fg\right)Gg''
\]
Nun ist $G-F$ eine Größe dritter Ordnung, also $1-\frac FG$ eine Größe zweiter Ordnung usw. Folglich ist auch $\frac FG.\frac{f''}{g''}-\frac{F''}{G''}.\frac fg$ von zweiter Ordnung und damit $Ff''-F''f$ von vierter Ordnung (sogar von fünfter Ordnung, wenn $\tau'$ in die Mitte zwischen $\tau$ und $\tau''$ fiele). Die rechten Seiten der zweiten, dritten und vierten Gleichung oben sind also von fünfter Ordnung, während auf der linken Seite jeweils sowohl der erste als auch der zweite Teil von dritter Ordnung ist. Für die erste Näherung kann man daher setzen
\[
\begin{array}{ll}
\text{aus 2)}&f'\delta'[\pi\pi'P']=-f''\delta''[\pi\pi''P']\\
\text{aus 4)}&f\delta[\pi\pi''P']=-f'\delta'[\pi'\pi''P']
\end{array}
\]
Was aus 3) folgt, ist mit diesen Ergebnissen identisch. Zur weiteren Abkürzung können wir hier auch $\frac fg=\frac{f''}{g''}$ setzen. Beide Größen unterscheiden sich von der Einheit nur um einen Betrag zweiter Ordnung und voneinander um einen Betrag derselben Ordnung (liegt $\tau'$ in der Mitte zwischen $\tau$ und $\tau''$, so ist die Differenz nur von dritter Ordnung). Da nun $g:-g':g''=\tau''-\tau':\tau''-\tau:\tau'-\tau$ gilt, erhält man
\[
\begin{aligned}
5)\qquad\delta&=\frac gf.\frac{f'}{g'}.\frac{\tau''-\tau}{\tau''-\tau'}.\frac{[\pi'\pi''P']}{[\pi\pi''P']}.\delta'\\
6)\qquad\delta''&=\frac{g''}{f''}.\frac{f'}{g'}.\frac{\tau''-\tau}{\tau'-\tau}.\frac{[\pi\pi'P']}{[\pi\pi''P']}.\delta'
\end{aligned}
\]

Diese Formeln ergeben $\delta$ und $\delta''$ aus $\delta'$ bis zur zweiten Ordnung einschließlich richtig, wenn $\tau'$ in der Mitte zwischen $\tau$ und $\tau''$ liegt, sonst ausschließlich der zweiten Ordnung. Im letzten Fall kann man $\frac{f'}{g'}=1$ setzen, da der Unterschied nur von zweiter Ordnung ist. Im ersten Fall hingegen lohnt es sich
\[
f'=g'+\tfrac43(f+f'+f'')\quad\text{oder}\quad\frac{f'}{g'}=1+\tfrac43\frac{f+f'+f''}{f'}
\]
zu setzen, was bald genauer bestimmt wird und um 1 Ordnung genauer ist. (Man sieht leicht, dass $f+f'+f''$ gleich dem Flächeninhalt des Dreiecks zwischen den drei Positionen
\SourceLeaf{162}%
\expandafter\def\csname HistoricalNoteText10\endcsname{Abhandlung über die leichteste Methode u.\ s.\ w.\ S.\ 45.}%
$p$, $p'$, $p''$ ist, nach einer bekannten Näherung also $=\frac34\times$ Flächeninhalt des Abschnitts zwischen der Sehne $pp''$ und dem Bogen.) Aus den obigen Formeln folgt ferner
\[
\frac{\delta''}{\delta}=\frac{[\pi\pi'P']}{[\pi'\pi''P']}.\frac{\tau''-\tau'}{\tau'-\tau}
\]
und wenn man für $\mathfrak{Z}$ die Ekliptik annimmt oder $B$, $B'$, $B''=0$ setzt, geht dies über in
\[
\frac{\delta''}{\delta}=\frac{\operatorname{tang}\beta\sin(\lambda'-L')-\operatorname{tang}\beta'\sin(\lambda-L')}{\operatorname{tang}\beta'\sin(\lambda''-L')-\operatorname{tang}\beta''\sin(\lambda'-L')}.\frac{\tau''-\tau'}{\tau'-\tau}
\]
also in die bekannte \textsc{Olbers}sche Formel\HistoricalNote{10}.

\begin{center}6.\end{center}
Nachdem wir aus den Formeln 2.\ 3.\ 4.\ bequeme Näherungen abgeleitet haben, gehen wir bei der ersten ähnlich vor. Bekanntlich gilt
\[
\begin{aligned}
\frac1r&=\frac1k(1-e\cos(v-\pi))\\
\frac1{r'}&=\frac1k(1-e\cos(v'-\pi))\\
\frac1{r''}&=\frac1k(1-e\cos(v''-\pi))
\end{aligned}
\]

Daraus folgt durch Multiplikation mit $\sin(v''-v')$, $\sin(v-v'')$, $\sin(v'-v)$ und anschließende Addition
\[
\begin{aligned}
\frac{f+f'+f''}{rr'r''}&=\frac1k(\sin(v''-v')+\sin(v-v'')+\sin(v'-v))\\
&=-\frac4k\sin\tfrac12(v''-v')\sin\tfrac12(v-v'')\sin\tfrac12(v'-v)
\end{aligned}
\]
oder
\[
\frac{f+f'+f''}{f'}=-\frac{2r'}k.\frac{\sin\tfrac12(v''-v')\sin\tfrac12(v'-v)}{\cos\tfrac12(v''-v)}
\]
Nach einem bekannten Lehrsatz der theoretischen Astronomie gilt
\[
\frac{\text{Überstrichene Fläche}}{\text{Mittlere Bewegung}}=\frac{a^{\frac32}.\BareRoot{k}}2
\]
Daraus folgt
\[
k=\frac{4gg''}{a^3(m'-m)(m''-m')}=\frac{4gg''}{A^3(M'-M)(M''-M')}
\]
\SourceLeaf{163}%
\expandafter\def\csname HistoricalNoteText11\endcsname{Hier wird tatsächlich subtrahiert, da $[\pi P\pi'']$, $[\pi P''\pi'']$ dasselbe Vorzeichen haben. Bei den Koeffizienten von $Ff''-F''f$ in den Gleichungen 2.\ 3.\ 4.\ werden die Teile hingegen tatsächlich addiert. Eine eingehendere Untersuchung würde hier zu weit führen.}%
Also gilt
\[
\begin{aligned}
\frac{f+f'+f''}{f'}&=-\frac{A^3(M'-M)(M''-M)}{2\cos^2\tfrac12(v''-v)}.\frac1{r'^3}.\frac{r'r'}{rr''}\\&\quad.\frac{r'r''\sin\tfrac12(v''-v')}{g}.\frac{rr'\sin\tfrac12(v'-v)}{g''}\\
\frac{f+f'+f''}{g'}&=-\frac{r'r''\sin\tfrac12(v''-v')}{g}.\frac{rr''\sin\tfrac12(v''-v)}{g'}.\frac{rr'\sin\tfrac12(v'-v)}{g''}\\&\quad.\frac{A^3(M'-M)(M''-M)}{2r'^3}.\frac{r'r'}{rr''}
\end{aligned}
\]
Man schließt daraus leicht, dass die Größen
\[
\frac1{\cos\tfrac12(v''-v)},\qquad\frac{r'r''\sin\tfrac12(v''-v')}{g},\qquad\frac{rr'\sin\tfrac12(v'-v)}{g'}
\]
und, falls entweder $\tau'$ in die Mitte zwischen $\tau$ und $\tau''$ fällt oder man die Abweichung der Bahn von $p$ von einem Kreis als Größe erster Ordnung betrachten kann, auch $\frac{r'r'}{rr''}$ von der Einheit nur um Größen zweiter Ordnung abweichen und man daher näherungsweise setzen darf
\[
\frac{f+f'+f''}{f'}=-\frac{A^3}{2r'^3}(M'-M)(M''-M')
\]
Entsprechend gilt näherungsweise
\[
\frac{F+F'+F''}{F'}=-\frac{A^3}{2R'^3}(M'-M)(M''-M')
\]

Die letzte Größe lässt sich auf Wunsch auch genau berechnen, da alle dafür nötigen Größen gegeben sind. Beide Größen sind von zweiter Ordnung und bis zur vierten ausschließlich bestimmt. — Wir haben also
\[
\begin{aligned}
F'(f+f'')-(F+F'')f'&=F'(f+f'+f'')-(F+F'+F'')f'\\
&=F'f'\tfrac12A^3(M'-M)(M''-M')\left(\frac1{R'^3}-\frac1{r'^3}\right)
\end{aligned}
\]
Diese Größe ist von vierter Ordnung und bis zur sechsten ausschließlich bestimmt. In Gleichung 1 oben ist die linke Seite von fünfter Ordnung. Auf der rechten Seite ist das erste Glied von sechster oder siebter Ordnung, denn $Ff''-F''f$ ist von vierter oder fünfter Ordnung und $D[\pi P\pi'']-D''[\pi P''\pi'']$ ist von zweiter\HistoricalNote{11} Ordnung. Das zweite Glied ist von fünfter Ordnung. Wir lassen daher das erste weg und erhalten
\[
7)\qquad\frac{[\pi\pi'\pi'']}{[\pi P'\pi'']}.\frac2{A^3(M'-M)(M''-M')}=\left(\frac1{R'^3}-\frac1{r'^3}\right)\frac R{\delta'}
\]
\SourceLeaf{164}%
mit einer Genauigkeit bis ausschließlich zur zweiten Ordnung, wenn $\tau'$ in die Mitte zwischen $\tau$ und $\tau''$ fällt. Sonst ist der Fehler nur von erster Ordnung. Diese Formel erhält die folgende Gestalt, wenn wir für $\mathfrak{Z}$ die Ekliptik nehmen. Sie ist der wichtigste Teil der ganzen Methode und ihre erste Grundlage.
\[
\begin{aligned}
\left\{1-\left(\frac{R'}{r'}\right)^3\right\}.\frac{R'}{\delta'}&=\frac2{A^3.(M'-M)(M''-M')}\\
&\quad.\frac{\operatorname{tg}\beta'\sin(\lambda''-\lambda)-\operatorname{tg}\beta\sin(\lambda''-\lambda')-\operatorname{tg}\beta''\sin(\lambda'-\lambda)}{\operatorname{tg}\beta\sin(L'-\lambda'')-\operatorname{tg}\beta''\sin(L'-\lambda)}
\end{aligned}
\]
$L'$ ist die Länge der Sonne $+180^0$.

Da die rechte Seite gegeben ist, sieht man, dass sich durch die Verbindung dieser Gleichung mit der folgenden
\[
\frac{\frac{R'}{\delta'}}{\frac{R'}{r'}}=\BareRoot{(1+\operatorname{tang}\beta'^2+\frac{R'R'}{\delta'\delta'}+2\frac{R'}{\delta'}\cos(\lambda'-L'))}
\]
$r'$ leicht finden lässt. Die indirekte Methode ist hier bei Weitem die bequemste. Mit wenigen Versuchen, für die sich leicht zweckmäßige Regeln angeben lassen, kommt man sehr schnell zum Ziel. Dabei erkennt man auch stets, ob es mehr als einen Wert für $r'$ und somit mehr als eine Bahn gibt, die die Beobachtungen darstellt. Das kann allerdings gelegentlich vorkommen. —

Weiter ist zu beachten, dass man die Längen hierbei eigentlich von einem festen Punkt statt von den beweglichen Äquinoktialpunkten aus zählen muss. Bei der Anwendung ist der Unterschied jedoch ohne Bedeutung. Drückt man die Zeit in Tagen aus, so gilt
\[
\log(M'-M)(M''-M')=\log(\tau'-\tau)+\log(\tau''-\tau')+6.4711352(-10)
\]
($M$ usw. müssen nämlich in Teilen des Radius statt in Grad ausgedrückt werden.)

Hat man $\delta'$ und $r'$, so lässt sich auch $\frac{f+f'+f''}{f'}$ bestimmen und damit auch $\delta$ und $\delta''$. Aus der Betrachtung der Formel 7) lassen sich noch weitere interessante Folgerungen ableiten, die hier übergangen werden müssen.

\begin{center}Zweiter Schritt.\\\emph{Näherungsweise Bestimmung der Bahnelemente.}\end{center}
\begin{center}7.\end{center}
Die mittlere Beobachtung zur Zeit $\tau'$ lassen wir nun ganz weg. An ihrer Stelle verwenden wir die Abstände $\delta$, $\delta''$, die im vorigen Schritt näherungsweise be%
\SourceLeaf{165}%
stimmt wurden. Daraus lassen sich nun offensichtlich die heliozentrischen Längen, Breiten und Abstände ableiten und daraus wiederum die Länge des $\AscendingNode$, die Neigung der Bahn und die Längen in der Bahn. Es bleibt also noch die Aufgabe,
\begin{center}\begin{tabular}{ll}
Aus zwei Längen in der Bahn $\quad.\quad.$&$v,\ v''$\\
\quad den Abständen von der Sonne $\quad.$&$r,\ r''$\\
\quad den zugehörigen Zeiten $\quad.\quad.\quad.$&$\tau,\ \tau''$
\end{tabular}\end{center}
die übrigen Elemente zu bestimmen, nämlich $a$, $e$, $\pi$ und die Epoche. Da die Beziehungen dieser Größen zu den gegebenen Größen transzendent sind, muss man wieder auf indirekte Methoden zurückgreifen. Wir betrachten hier drei davon.

\par\medskip\emph{Erste Methode.}\par
Wenn $\pi$ als gegeben vorausgesetzt wird.

Man setze
\[
\frac{r''+r}{r''-r}\operatorname{tang}\tfrac12(v''-v)=\operatorname{tang}\zeta
\]
so ist
\[
e=\frac{\cos\zeta}{\cos\tfrac12(v''-v)\cos[\pi-\tfrac12(v+v'')-\zeta]}
\]
Am besten berechnet man $k$ anschließend auf zwei Arten
\[
k=r[1-e\cos(v-\pi)]=r''[1-e\cos(v''-\pi)]
\]
Das dient zugleich zur Kontrolle der Rechnung. Setzt man $e=\sin\varphi$, so ist $a=\frac{k}{\cos\varphi^2}$. Aus den wahren Anomalien lassen sich dann die exzentrischen und mittleren Anomalien und Längen mit den üblichen Methoden oder bequemer mit indirekten Methoden berechnen. Daraus und aus der durch $a$ gegebenen mittleren Bewegung erhält man zwei Bestimmungen der mittleren Länge für eine beliebige Epoche. Stimmen beide überein, so hat man den richtigen Wert von $\pi$ gefunden. Andernfalls muss man die Rechnung mit einem etwas veränderten Wert von $\pi$ wiederholen und den wahren Wert durch Interpolation bestimmen. ‘Dabei empfiehlt es sich, die übrigen Elemente \emph{nicht} durch Interpolation zu bestimmen, sondern sie aus den korrigierten Werten von $\pi$ neu zu berechnen und erst aufzuhören, wenn die beiden Werte für die Epoche vollkommen übereinstimmen.’
\SourceLeaf{166}%
\par\medskip\emph{Zweite Methode.}\par
Wenn $e$ vorausgesetzt wird.

Die Rechnung ist genau dieselbe. Da man hier jedoch $\pi$ aus der Gleichung
\[
\cos[\pi-\tfrac12(v+v'')-\zeta]=\frac{\cos\zeta}{e\cos\tfrac12(v''-v)}
\]
bestimmen muss, braucht man bereits eine ungefähre Kenntnis des wahren Werts, weil dem Kosinus zwei verschiedene Werte entsprechen.

Methode I ist II vorzuziehen. Beide Methoden sind überhaupt nur dann zweckmäßig, wenn der durchlaufene Bogen schon sehr groß ist und man die Elemente schon beinahe kennt. Für die ersten Näherungen aus einer kurzen Beobachtungsreihe verwendet man stets die folgende Methode.

\Needspace{6\baselineskip}\par\medskip\emph{Dritte Methode.}\par
Wenn $k$ vorausgesetzt wird.

Man hat hier
\[
\begin{aligned}
\frac{\frac1{r''}-\frac1r}{2\sin\tfrac12(v''-v)}&=\frac ek\sin(\tfrac12(v+v'')-\pi)\\
\frac{\frac2k-\frac1{r''}-\frac1r}{2\cos\tfrac12(v''-v)}&=\frac ek\cos(\tfrac12(v+v'')-\pi)
\end{aligned}
\]
Durch Division erhält man also $\operatorname{tang}[\tfrac12(v+v'')-\pi]$, daraus $\pi$ und dann aus einer der beiden Gleichungen $e$. Alles Weitere geschieht genau wie bei den vorhergehenden Methoden.

Der Vorteil der dritten Methode liegt darin, dass man für $k$ sofort einen sehr guten Näherungswert finden kann, wenn der Bogen $v''-v$ nicht zu groß ist. Für den Ausschnitt zwischen den beiden Radiusvektoren gilt nämlich
\[
g'=\frac{a^{\frac32}\BareRoot{k}}2(m''-m)=\tfrac12A^{\frac32}(M''-M)\BareRoot{k}
\]
Nun ist $2g'=\int rr\,\mathrm{d}w$ von $w=v$ bis $w=v''$.

Nach der bekannten Formel von \textsc{Cotes} zur näherungsweisen Integration ist nun $\int\varphi w.\mathrm{d}w$ von $w=v$ bis $w=v''$
\[
=(\tfrac12\varphi v+\tfrac12\varphi v'')(v''-v)
\]
\SourceLeaf{167}%
und noch genauer
\[
=(\tfrac16\varphi v+\tfrac23\varphi\tfrac12(v+v'')+\tfrac16\varphi v'')(v''-v)
\]
noch genauer
\[
=(\tfrac18\varphi v+\tfrac38\varphi(\tfrac23v+\tfrac13v'')+\tfrac38\varphi(\tfrac13v+\tfrac23v'')+\tfrac18\varphi v'')(v''-v)\quad\text{usw.}
\]
Es reicht jedoch immer aus, bei den ersten beiden Formeln zu bleiben.

Nach der ersten hat man also
\[
2g'=\tfrac12(rr+r''r'')(v''-v)\quad\text{und}\quad\BareRoot{k}=\frac{\tfrac12(rr+r''r'')}{A^{\frac32}}.\frac{v''-v}{M''-M}
\]
Gewöhnlich setzt man $A$ $=1$. Drückt man $v''-v$ und $M''-M$ in Bogensekunden aus, so ist $\log(M''-M)=\log(\tau''-\tau)+3.5500073$. Zur Abkürzung der Rechnung setzt man $\frac{r'}r=\operatorname{tang}(45^0\pm\psi)$, woraus $\tfrac12(rr+r''r'')=\frac1{\cos2\psi}$ folgt.

Für die zweite Integrationsformel setzt man den Radiusvektor, der zur Länge $\tfrac12(v''+v)$ gehört, $=r^*$. Dann ist
\[
\frac1{r^*}=\tfrac12\left(\frac1r+\frac1{r''}\right)+\left[\tfrac12\left(\frac1r+\frac1{r''}\right)-\frac1k\right]\frac{2\sin\tfrac14(v''-v)^2}{\cos\tfrac12(v''-v)}
\]
Daraus kann man $r^*$ mithilfe des ersten Werts von $k$ bestimmen. Dann gilt
\[
2g'=\tfrac16(rr+r''r'')+\tfrac23r^*r^*
\]
und der neue Wert von $k$ ist damit
\[
=k\left\{1+\frac{\tfrac23(r^*r^*-\tfrac12(rr+r''r''))}{rr+r''r''}\right\}^2
\]

In der Praxis ist es gewöhnlich genau genug und bequemer, den Logarithmus des neuen Werts von $k$ zu bestimmen, indem man den Logarithmus des ersten Werts um $\tfrac43\log\frac{r^*r^*}{\tfrac12(rr+r''r'')}$ erhöht. Bestimmt man mit diesem neuen Wert von $k$ den Wert von $\frac1{r^*}$ nach der obigen Formel nochmals genauer und korrigiert dann den Wert von $k$ erneut, so stimmen die beiden zuletzt erhaltenen Bestimmungen der Epoche fast immer so gut überein, dass überhaupt keine neue Annahme nötig ist. Bei $\text{\AstroSymbols ⚳}$ und $\text{\AstroSymbols ⚴}$ stimmten sie, obwohl $\tau''-\tau$ 41 und 42 Tage betrug, immer bis auf ein paar Hundertstelsekunden überein.
\SourceLeaf{168}%
\begin{center}\emph{Verbesserungsverfahren.}\end{center}
\begin{center}8.\end{center}
Berechnet man aus den mit den vorhergehenden Methoden gefundenen Näherungswerten der Elemente die Position zur Zeit $\tau'$ und stimmt diese mit der Beobachtung überein, so ist die Arbeit abgeschlossen. Die Übereinstimmung wird gewöhnlich sehr gut sein (bei meinen Rechnungen betrug der Unterschied oft nur wenige Bogensekunden), aber selten vollkommen. Das liegt zum einen daran, dass die zugrunde gelegten Annahmen zum Teil nur näherungsweise gelten, zum anderen \emph{daran, dass selbst die dabei verwendeten Positionen der Sonne keiner Ellipse folgen, sondern die kleinen Störungen mit einschließen}. Man könnte zwar die Werte der oben weggelassenen kleinen Größen höherer Ordnung aus den genäherten Elementen sehr genau bestimmen und damit die obigen Formeln und die Werte von $\delta$, $\delta''$ verbessern. Ich meine jedoch, dass diese Rechnung viel mühsamer wäre als eine der folgenden Methoden.

Das einfachste Verfahren für eine erste Verbesserung ist das folgende. Darauf kam ich erst bei der Arbeit an $\text{\AstroSymbols ⚴}$ und wandte es dabei mit größtem Erfolg an, da die Zeitabstände noch kurz waren.

Angenommen, für die genäherten Elemente, \emph{die auf die oben beschriebene Weise gefunden wurden}, sei die berechnete Position zur Zeit $\tau$ in der Länge $=\lambda'+\mathfrak{L}$ und in der Breite $=\beta'+\mathfrak{B}$, während die beobachteten Werte $\lambda'$ und $\beta'$ sind. Alle kleinen Ungenauigkeiten in den Annahmen verschwören sich also, sodass die Länge um $\mathfrak{L}$ und die Breite um $\mathfrak{B}$ zu groß ausfällt. Dann berechnet man die Bahn ganz von neuem und \emph{auf genau dieselbe Weise}, indem man die Beobachtungen
\[
\begin{array}{ccc}
\lambda,&\lambda'-\mathfrak{L},&\lambda''\\
\beta,&\beta'-\mathfrak{B},&\beta''
\end{array}
\]
zugrunde legt. Die aus den neuen Elementen berechnete Position wird sich dann von $\lambda'$, $\beta'$ so wenig unterscheiden (nach meiner Erfahrung nur um Bruchteile einer Bogensekunde), dass keine weitere Verbesserung nötig ist.

\begin{center}9.\end{center}
Das eben beschriebene Verfahren gilt nur, wenn man die Bahnbestimmung auf \emph{dieselben} Beobachtungen stützen will, die man für die erste An%
\SourceLeaf{169}%
\expandafter\def\csname HistoricalNoteText12\endcsname{Siehe \textit{Theoria motus corporum coelest.}, Art.\ 83.}%
näherung verwendet hat. Will man aber danach die Elemente anhand ganz oder teilweise anderer Beobachtungen verbessern, so habe ich nach verschiedenen anderen Versuchen die folgenden beiden Methoden am brauchbarsten gefunden.

I. Man berechnet aus den zwei äußeren geozentrischen Positionen die heliozentrischen Positionen nach 3 Hypothesen. Zuerst setzt man die genäherten Abstände für diese Beobachtungen voraus. Dann ändert man zunächst den einen und danach den anderen Abstand ein wenig. Mit den unter den 3 Hypothesen gefundenen Elementen berechnet man jeweils die Position für die mittlere Beobachtung und vergleicht sie mit der beobachteten Position. Durch Interpolation erhält man dann die korrigierten Abstände und auf Wunsch auch die korrigierten Elemente. Es ist jedoch besser, die Mühe auf sich zu nehmen, die Elemente gesondert aus den neuen Abständen zu berechnen, besonders wenn sich die Elemente noch sehr stark ändern.

II\textsuperscript{a}. Man geht genau so vor, verwendet aber anstelle der genäherten Abstände für die äußeren Beobachtungen die genäherten Bestimmungen der Neigung und des $\AscendingNode$ und ändert jede von beiden ein wenig.

II\textsuperscript{b}. Man berechnet aus allen drei geozentrischen Positionen die heliozentrischen, teils mit den genäherten, teils mit etwas veränderten Bestimmungen der Neigung und des $\AscendingNode$. Aus den zwei äußeren heliozentrischen Positionen berechnet man die Elemente und aus diesen die mittlere heliozentrische Position. Diese vergleicht man mit der Position, die man aus der beobachteten geozentrischen Position abgeleitet hat. Dann bestimmt man durch Interpolation die verbesserte Neigung und den verbesserten $\AscendingNode$ usw.

In II\textsuperscript{b} könnte man auch mit bekannten Formeln aus den \emph{drei heliozentrischen} Positionen die Ellipse bestimmen, ohne die Zeiten heranzuziehen. Aus den Abmessungen der Ellipse ließen sich dann die 2 Zeitabstände berechnen und mit den wahren vergleichen. Anschließend könnte man wie zuvor die korrigierte Neigung und den korrigierten $\AscendingNode$ durch Interpolation bestimmen. Nach meiner Erfahrung musste ich dieses Verfahren jedoch verwerfen. Auf diesem Weg würde man erst nach wiederholten Rechenschritten mit weit größerer Mühe zu einer genauen Darstellung der Beobachtungen gelangen. Die Gründe dafür hier im Einzelnen zu untersuchen, würde zu weit führen\HistoricalNote{12}. Ich bemerke daher nur, dass man auf diese Weise die zweiten Dif%
\SourceLeaf{170}%
ferenziale wieder einführt, die man gerade durch die in den Artikeln 5.\ und 6.\ verwendeten Kunstgriffe beseitigt hat. Diese empfindlichen zweiten Differenziale können durch eine geringe Änderung der Neigung und des $\AscendingNode$ \emph{ganz enorm} entstellt werden, besonders wenn die Exzentrizität gering ist. Schon eine Änderung um ein paar Bogenminuten beim $\AscendingNode$ oder bei der Neigung kann leicht eine Ellipse hervorbringen, die mit der vorherigen fast keine Ähnlichkeit mehr hat. Der Interpolation kann man dann verständlicherweise nicht mehr trauen. Bei unseren Methoden tritt das nicht ein, da sie stets auf nur zwei Beobachtungen beruhen. \emph{Sapienti sat.} — Nach meinen wiederholten Erfahrungen halte ich Methode I für die zweckmäßigste und allgemeinste.

Alle diese Methoden gelten nur, solange der Bogen noch mäßig groß ist. Liegen bereits Beobachtungen aus 1 oder mehreren Jahren vor, so braucht man wieder andere Methoden, auf die ich hier nicht ausführlich eingehen kann. In diesem Fall ist es im Allgemeinen nicht ratsam, die Elemente aus drei vollständigen Beobachtungen zu bestimmen. Weit angemessener ist es, vier Längen und zwei Breiten zu verwenden. — Umfassen die Beobachtungen schon \emph{mehrere} Jahre und sind die Elemente bereits \emph{bis auf Kleinigkeiten} bestimmt, so halte ich Differentialkorrekturen für das beste Mittel. Dabei kann man eine beliebige Zahl von Beobachtungen zugrunde legen.
\begin{center}\rule{2.65cm}{0.4pt}\end{center}
